\chapter{Recent Transmission Compatibility}\label{ch:epilink}

\section{Introduction}\label{sec:ch3_introduction}

Section~\ref{sec:ch2_common_approaches} distinguishes reconstruction of who infected whom from the narrower surveillance question of whether sampled infections are compatible with recent linkage. Genetically similar infections can arise through different transmission histories, while differences between sampling dates also reflect delays to observation. EpiLink combines these two observations under explicit transmission scenarios to identify pairs and groups worth investigating when labelled transmission pairs are unavailable. This chapter develops and evaluates that approach. It is adapted from the preprint by \citet{Arthur2026}, which was under review at \textit{PLOS Computational Biology} during preparation of this thesis.

The scenarios evaluated here are direct transmission between sampled infections and infection from a shared source. EpiLink compares observed genetic distances and sampling-time differences with simulated values under these scenarios. The simulations represent variation in infection timing, testing delay, and mutation accumulation under the chosen model assumptions. The resulting scores express relative compatibility and provide weights for constructing graphs and identifying clusters. They are not transmission probabilities and do not identify who infected whom.

The evaluation uses synthetic SARS-CoV-2 data with known links to assess how well EpiLink distinguishes linked pairs and recovers clusters, how performance changes under different model assumptions, and how stable the clusters remain as data accumulate. Comparisons include logistic regression trained on labelled pairs and phylogenetic clustering with TreeCluster. An application to the Boston epidemic examines whether the inferred clusters recover documented outbreaks \citep{Lemieux2021}.

\section{Materials and Methods}\label{sec:ch3_methods}

\subsection{EpiLink compatibility model}\label{sec:ch3_epilink_model}

EpiLink is a pairwise compatibility model that assesses whether the observed consensus genetic distance $g_{ij}$ and sampling-time difference $t_{ij}$ between two sampled cases $(i,j)$ are consistent with specified transmission relationships (Figure~\ref{fig:epilink_model}). It represents variation in infection timing, testing delay, and mutation accumulation under chosen parameter values, without reconstructing a unique transmission tree. Sensitivity analyses examine alternative parameter values separately; these simulations do not integrate over a posterior distribution of parameter uncertainty.

\begin{figure}[htbp]
    \centering
    \includegraphics[width=\textwidth]{epilink_schematic}
    \caption[Overview of the EpiLink compatibility model]{\textbf{Overview of the EpiLink compatibility model.} The evaluated target set contains direct transmission and infection from a shared source, with no additional unsampled intermediates. For a sampled pair $(i,j)$, observed sampling-time difference $t_{ij}$ and consensus genetic distance $g_{ij}$ are compared with simulated values under each scenario $s$. The empirical distribution functions $\widehat F_{T,s}$ and $\widehat F_{G,s}$ give the proportions of simulated values at or below each observation. Their transformed compatibility scores are multiplied within each scenario, $C_s(i,j)=C_{T,s}(i,j)C_{G,s}(i,j)$, then summed over the target set $\mathcal{S}_\star$.}\label{fig:epilink_model}
\end{figure}

The analyses in this chapter evaluate two scenarios: direct transmission between the sampled cases, and co-primary infection in which both sampled cases were infected by the same unsampled source. We write this target set as

\[
	\mathcal{S}_\star =
	\{H_{\mathrm{AD}}(0),\;H_{\mathrm{CA}}(0,0)\}.
\]

Here, AD denotes an ancestor-descendant relationship and CA a common-ancestor relationship. The zeros indicate that no additional unsampled intermediates separate the sampled cases from the direct infector or shared source; the shared source itself is not counted as an intermediate. The general scenario space allowing up to $M$ intermediates is defined in Appendix Section~\ref{sec:app_ch3_latent_scenario_space}. All evaluations here use $M=0$.

For each scenario $s\in\mathcal{S}_\star$, EpiLink generates Monte Carlo draws from a shared process-based model with Gamma-distributed latent ($E$), presymptomatic infectious ($P$), and symptomatic infectious ($I$) stage durations, together with a Gamma-distributed delay from symptom onset to testing or sampling \citep{Hart2021,McAloon2020,Hart2022}. Writing the incubation period as $\tau_{\mathrm{inc}}=y_E+y_P$ and the generation time as $\tau_{\mathrm{gen}}=y_E+y_{\mathrm{toit}}$, where $y_{\mathrm{toit}}$ is the time from onset of infectiousness to transmission, the total observation interval from exposure to testing is $\tau_{\mathrm{obs}}=\tau_{\mathrm{inc}}+x_{\mathrm{test}}$.

For each scenario, these simulations produce a distribution of possible differences between testing dates, $T_s$. The same unobserved history defines an effective branch length $B_s$, representing evolutionary time between the sampled genomes. For ancestor-descendant scenarios, this combines the generation intervals along the chain with the difference between the two exposure-to-observation intervals, with negative values set to zero. For common-ancestor scenarios, it adds the generation intervals along both branches and both exposure-to-observation intervals. The full derivation of $T_s$ and $B_s$ is given in Appendix Section~\ref{sec:app_ch3_model_derivation}.

Branch-length draws were mapped to expected genetic distances using a per-genome daily substitution rate $\lambda=\nu L/365$, where $\nu$ is the per-site per-year substitution rate and $L$ is the genome length. Under the deterministic genetic model, $G_s=\lambda B_s$; under the stochastic genetic model, $G_s\mid B_s,\lambda\sim\mathrm{Poisson}(\lambda B_s)$ \citep{Jombart2010}. This choice is separate from the molecular clock: a strict clock fixes the substitution rate, while a relaxed uncorrelated log-normal clock draws $\nu$ from a log-normal distribution \citep{Drummond2006}.

Compatibility is calculated by locating each observation within its simulated distribution. Let $X$ denote either the temporal quantity $T$ or the genetic quantity $G$, and let $X_s^{(r)}$ be Monte Carlo draw $r$ under scenario $s$. With $N_s$ draws, the empirical cumulative distribution function is

\[
    \widehat F_{X,s}(x)
    = \frac{1}{N_s}\sum_{r=1}^{N_s} 
    \mathbf{1}\!\left\{X_s^{(r)}\le x\right\}, 
    \qquad X\in\{T,G\}.
\]

The indicator $\mathbf{1}\{\cdot\}$ equals one when its condition holds and zero otherwise, so $\widehat F_{X,s}(x)$ is the proportion of simulated values at or below $x$. It is an empirical percentile expressed on a zero-to-one scale, not a significance $p$-value. The scalar compatibility transformation is

\[
    c_{X,s}(x)=1-2\left|\widehat F_{X,s}(x)-\tfrac{1}{2}\right|.
\]

It equals one when the empirical percentile is one half and falls towards zero as that percentile approaches either endpoint. For example, percentiles 0.25 and 0.75 both give compatibility 0.5. With discrete genetic draws, ties can prevent a score of exactly one; Appendix Section~\ref{sec:app_ch3_compatibility_score} explains this point.

For the observed pair, the temporal and genetic components are

\[
    C_{T,s}(i,j)=c_{T,s}(t_{ij}),
    \qquad C_{G,s}(i,j)=c_{G,s}(g_{ij}).
\]

We multiply these components within each scenario and sum the resulting scenario scores:

\[
    C_s(i,j)=C_{T,s}(i,j)C_{G,s}(i,j),
    \qquad
    C_{\mathcal{S}_\star}(i,j)=\sum_{s\in\mathcal{S}_\star}C_s(i,j).
\]

\subsection{Interpretation, scope, and assumptions}\label{sec:ch3_epilink_interpretation}

The score $C_{\mathcal{S}_\star}(i,j)$ is a compatibility weight for a specified set of recent linkage hypotheses. It measures how typical the observed testing-time gap and consensus-genome distance are compared with values simulated under the selected scenarios $\mathcal{S}_\star$ and model parameters. A high score indicates that a pair lies near the central region of the time and genetic distributions generated by at least one selected scenario. It does not identify the infector, infer transmission direction, distinguish direct transmission from co-primary infection without inspecting the scenario-level components, or estimate a calibrated posterior probability of transmission. Conversely, a low score indicates a poor match to the selected scenarios and assumptions; it does not establish that epidemiological linkage is absent.

The absolute scale of $C_{\mathcal{S}_\star}(i,j)$ depends on the number and definition of the selected scenarios because the final score sums scenario-level compatibilities. Scores can therefore exceed 1 when more than one scenario aligns well with the same observed pair. In the default configuration used here, the target is restricted to direct transmission and co-primary infection. Adding scenarios with unsampled intermediates to $\mathcal{S}_\star$ would include longer transmission chains, so a high score would be less specific to direct transmission or infection from a shared source. The main assumptions that determine this interpretation are summarised in Table~\ref{tab:epilink_assumptions}.

\begin{table}[htbp]
    \centering
    \caption[Interpretation and assumptions of the EpiLink compatibility score]{Interpretation and assumptions of the EpiLink compatibility score. The table summarises the assumptions that define the scope of the score and the corresponding consequences for interpretation.}\label{tab:epilink_assumptions}
    \begin{thesistablebody}{@{}P{0.28\textwidth}P{0.62\textwidth}@{}}
    \toprule
    \textbf{Scope or assumption} & \textbf{Consequence for interpretation} \\
    \midrule
    Unsampled intermediates and selected scenarios &
    The evaluated target set contains direct transmission and co-primary infection, with no additional unsampled intermediates ($M=0$). Longer chains are available in the general formulation in Appendix~\ref{app:epilink}, but are not evaluated here. \\

    Shared process-based model &
    Stage-duration, infectiousness, and testing-delay distributions are shared across cases within a run. Individual variation is represented only through these distributions or through alternative parameter settings. \\

    Transmission timing relative to symptom onset &
    Transmission timing is generated from an $E/P/I$ infectiousness profile relative to symptom onset, with a higher transmission rate before symptom onset and a lower rate afterwards \citep{Hart2021,He2020,Ferretti2020}. \\

    Testing date as observation time &
    The only observed time is the sample collection or testing date. Infection times are latent, and symptom-to-testing delays are modelled as independent draws, so shared or context-specific testing behaviour can affect compatibility. \\

    Independent generation intervals &
    Each generation interval comes from an independent draw from the same model once parameters are fixed. Correlations specific to households, workplaces, care homes, or hospitals are not explicitly modelled. \\

    Consensus-genome representation &
    The genetic input is consensus distance. Within-host diversity, minority variants, and explicit within-host coalescent time are not modelled, limiting resolution when consensus diversity is low \citep{Lythgoe2021}. \\

    Genetic distance given branch length and substitution rate &
    Genetic distance is treated as deterministic or Poisson-distributed conditional on effective branch length and substitution rate. Incorrect assumptions about the substitution rate, variation in that rate, or random variation in mutation counts can change compatibility scores. \\

    Combining time and genetic scores by multiplication &
    A pair receives a high scenario-level score only when it is compatible in both time and genetics under the same unobserved scenario. Adding the scores or taking their maximum would combine the time and genetic information differently. \\

    Compatibility rather than probability &
    Scores rank alignment with the chosen assumptions. They are not calibrated false-positive rates, posterior transmission probabilities, evidence of transmission direction, or proof of direct transmission. \\

    Using scores to construct graphs &
    Pairwise scores are used as edge weights before sparsification and Leiden community detection. Resulting clusters depend on graph-construction and resolution choices. They represent groups of sampled infections compatible with recent transmission, without establishing confirmed transmission events. \\
    \bottomrule
    \end{thesistablebody}
\end{table}

\subsection{Synthetic benchmark dataset}\label{sec:ch3_synthetic_benchmark}

We evaluated EpiLink using synthetic outbreak datasets from the Scotland Coronavirus transmission Model (SCoVMod)\nomenclature{SCoVMod}{Scotland Coronavirus transmission Model}, a spatially explicit agent-based model that simulates transmission across Scotland using demographic structure, mobility patterns, and time-varying contact behaviour \citep{Banks2022}. We constructed a directed transmission graph from its infection-history and transmission-event outputs. For cases with more than one recorded infector, we retained the earliest recorded infector and selected a group of approximately 5,000 cases connected when transmission direction was ignored. We then constructed a rooted directed transmission tree, in which each case except the root has one retained infector. The algorithm selects the tree with the largest total edge weight (a maximum spanning arborescence); negative simulation times were used as weights to favour earlier recorded infectors. The resulting tree comprised 4,990 nodes and 4,989 directed edges.

Transmission heterogeneity in the underlying SCoVMod outbreak was summarised by fitting a negative-binomial offspring distribution to the full transmission graph. This gave an estimated dispersion parameter of $\hat{k}=0.233$ (95\% confidence interval (CI)\nomenclature{CI}{confidence interval}: 0.232--0.235) and a mean of approximately 1.00 onward infections per individual (95\% CI: 0.992--1.006). Approximately 15.6\% of individuals accounted for 80\% of onward transmission, consistent with the superspreading dynamics motivating this work. We then used the EpiLink process-based model to simulate dates and genetic distances while keeping the selected transmission tree fixed. % chktex 13

\subsection{Baseline parameters and sensitivity analysis}\label{sec:ch3_parameters_sensitivity}

Key baseline evaluation parameters (Table~\ref{tab:epilink_parameters}) were drawn from published SARS-CoV-2 estimates \citep{Hart2021,Duchene2020,vanDorp2020}. Other $E/P/I$ infectiousness profile parameters were kept fixed at their fitted values, as described in Appendix Section~\ref{sec:app_ch3_hart_model}. We reduced or increased the incubation mean and coefficient of variation (CV)\nomenclature{CV}{coefficient of variation}, testing-delay mean and CV, and substitution rate by 25\%. For clock relaxation, we tested zero rate variation and a 25\% increase. This gave 12 sensitivity-analysis scenarios.

The synthetic sequences contained 5,000 sites. Separately, the model used $L=29{,}903$ to convert the per-site substitution rate to a per-genome rate during both data generation and inference. Thus, the simulated sequence length and the length used to scale mutation counts had different roles. The Boston alignment also contained 29,903 nucleotides.

\begin{table}[htbp]
    \centering
    \caption[Baseline evaluation parameters and perturbation design]{Baseline evaluation parameters and perturbation design.}\label{tab:epilink_parameters}
    \begin{thesistablebody}{@{}llcc@{}}
    \toprule
    \textbf{Parameter} & \textbf{Baseline value} & \textbf{Lower multiplier} & \textbf{Upper multiplier} \\
    \midrule
    Incubation mean & 5.505 days & 0.75 & 1.25 \\
    Incubation CV & 0.415 & 0.75 & 1.25 \\
    Testing-delay mean & 1.0 day & 0.75 & 1.25 \\
    Testing-delay CV & 1.0 & 0.75 & 1.25 \\
    Substitution rate & $10^{-3}$ sub/site/year & 0.75 & 1.25 \\
    Clock relaxation & 0.33 & 0.0 & 1.25 \\
    \midrule
    \multicolumn{4}{l}{\textit{Fixed simulation parameters:}} \\
    Simulated sequence length & 5,000 sites & --- & --- \\
    Length for rate scaling ($L$) & 29,903 sites & --- & --- \\
    Monte Carlo draws & 10,000 & --- & --- \\
    Target subset & $\{H_{\mathrm{AD}}(0),H_{\mathrm{CA}}(0,0)\}$ & --- & --- \\
    Sparsification threshold & $10^{-4}$ & --- & --- \\
    \bottomrule
    \end{thesistablebody}
\end{table}

Each perturbation scenario was evaluated under two inference conditions. In the matched condition, inference used the parameter values used to generate the data. In the mismatched condition, inference retained the baseline values. Comparing these conditions shows how performance changes when parameters are correctly specified and when they are not. Matching parameter values does not make the data-generation and inference formulations identical.

\subsection{Model variants and comparators}\label{sec:ch3_epilink_variants}

The comparisons separate two questions: how additional variation in mutation counts changes the information available in the data, and whether explicitly modelling that variation changes inference on the same data. Deterministic genetic inputs provide a controlled baseline using rounded expected counts; stochastic inputs add Poisson variation. Comparing the deterministic and stochastic inference formulations on each input type then isolates the choice made by the scoring model. These are diagnostic comparisons, rather than alternative ways to generate an observed genome.

We denote the deterministic-genetics and stochastic-genetics inference formulations by ED\nomenclature{ED}{EpiLink with deterministic-genetics inference} and ES\nomenclature{ES}{EpiLink with stochastic-genetics inference}, respectively. ES models mutation counts using Poisson variation around the expected number of substitutions. The three-letter benchmark labels append D or S to identify how the synthetic input data were generated: ESD\nomenclature{ESD}{EpiLink stochastic-genetics inference with deterministic synthetic input} and ESS\nomenclature{ESS (EpiLink)}{EpiLink stochastic-genetics inference with stochastic synthetic input} therefore use the same ES formulation with different types of synthetic input. For the Boston and Scottish observations, we use ES without an input-generation suffix. This choice represents realised mutation counts as variable around their expectation; it does not depend on being able to choose how the observed genomes were generated. The deterministic formulation remains a comparison for assessing the contribution of that additional variation.

EpiLink was supplied with synthetic pairwise sampling-time differences and genetic distances. Genetic distances were generated using either rounded expected mutation counts (D) or counts drawn with Poisson variation (S). This yields four inference/input combinations: EDD\nomenclature{EDD}{EpiLink deterministic-genetics inference with deterministic synthetic input}, EDS\nomenclature{EDS}{EpiLink deterministic-genetics inference with stochastic synthetic input}, ESD, and ESS (Table~\ref{tab:epilink_variants}). The D/S suffix describes genetic input generation, not the time data or a choice available when analysing observed genomes.

Two logistic regression models, LD\nomenclature{LD}{logistic regression with deterministic synthetic input} and LS\nomenclature{LS}{logistic regression with stochastic synthetic input}, were fitted as supervised reference comparators using the same pairwise features supplied to EpiLink. Both were trained on a random 10\% subsample of all case pairs, with at least one pair from each class. Pairwise discrimination was evaluated on all pairs, including the training pairs, and their scores also contributed to the graphs used for clustering. The comparison therefore describes performance on the full benchmark dataset, rather than a held-out evaluation of the logistic models.

\begin{table}[htbp]
    \centering
    \caption[EpiLink variants and supervised comparators used in the synthetic benchmark]{EpiLink variants and supervised comparators used in the synthetic benchmark.}\label{tab:epilink_variants}
    \begin{thesistablebody}{@{}P{0.12\textwidth}P{0.24\textwidth}P{0.24\textwidth}P{0.30\textwidth}@{}}
    \toprule
    \textbf{Label} & \textbf{Inference model} & \textbf{Input data-generating model} & \textbf{Interpretation} \\
    \midrule
    EDD & EpiLink deterministic genetics & Deterministic input & Uses expected mutation counts in generation and inference, with rounding during generation. \\
    EDS & EpiLink deterministic genetics & Stochastic input & Tests deterministic inference under noisy mutation counts. \\
    ESD & EpiLink stochastic genetics & Deterministic input & Stochastic inference applied to rounded expected-count inputs. \\
    ESS & EpiLink stochastic genetics & Stochastic input & Most directly matched to stochastic mutation-count input. \\
    LD & Logistic regression & Deterministic input & Supervised benchmark with access to transmission labels. \\
    LS & Logistic regression & Stochastic input & Supervised benchmark under noisy mutation-count input. \\
    \bottomrule
    \end{thesistablebody}
\end{table}

Graph construction and clustering used these pairwise outputs as edge weights: EpiLink compatibility scores for the four EpiLink variants and logistic probabilities for LD and LS. TreeCluster was used as a comparator that groups sequences using a dated phylogenetic tree \citep{Balaban2019}. A dated phylogeny was inferred from the synthetic benchmark data using TreeTime \citep{Sagulenko2018}, and TreeCluster was applied across distance thresholds expressed as calendar days divided by 365. TreeCluster assigns each sequence to one cluster without producing pairwise linkage scores, so pairwise ranking metrics such as average precision (AP)\nomenclature{AP}{average precision} were not calculated for it. % chktex 13

\subsection{Graph construction, community detection, and metrics}\label{sec:ch3_graph_clustering_metrics}

Total graph weight is the sum of the pairwise scores assigned to its edges. Weight retention is the fraction of that sum remaining after low-weight edges are removed; it differs from the fraction of edges retained. We tested edge-weight thresholds in $\{10^{-4},10^{-3},10^{-2},10^{-1}\}$ to reduce graph density while retaining at least 99.5\% of total edge weight. Among qualifying thresholds, the selection rule prioritised the greatest retained weight, then fewer retained edges, then the lower threshold if both measures were tied. This rule selected $10^{-4}$ for every model. Because retaining weight takes priority in the rule, agreement between selected thresholds is not independent evidence for a universal cut-off. We used the same threshold for Boston as a fixed operational choice, rather than selecting it using the known exposure groups. Weight retention describes the numerical effect of sparsification; it does not establish that all epidemiologically important connections remain. Appendix Section~\ref{sec:app_ch3_evaluation_definitions} gives the calculation.

For all synthetic experiments, weighted graphs were constructed from pairwise compatibility scores or logistic probabilities after excluding edges with weights below the sparsification threshold ($10^{-4}$). Community detection was performed using the Leiden algorithm \citep{Traag2019} with a modularity-based objective. We tested resolution values $\gamma$ from 0.1 to 1.0 in steps of 0.1, with 10 random restarts at each value. For each model, we selected the value giving the highest BCubed F1\nomenclature{F1}{harmonic mean of precision and recall}, the harmonic mean of BCubed precision and recall. Reported benchmark F1 values therefore use resolutions selected against the reference neighbourhoods.

Clustering performance was assessed using extended BCubed precision and recall \citep{Amigo2008}. Precision describes how consistently cases placed together agree with the reference memberships; recall describes how much of the reference association is retained by the inferred grouping. The extended measure allows a case to belong to more than one reference group. The reference neighbourhood for each transmission-tree case $u$ contains $u$ and its direct infectees. F1 is the harmonic mean of precision and recall. Appendix Section~\ref{sec:app_ch3_evaluation_definitions} states the membership and averaging conventions.

The reference builder was corrected and validated against all 4,990 cases before computing the BCubed results reported here. The completed correction and reproduction checks are recorded in Appendix Section~\ref{sec:app_ch3_reference_audit}.

Pairwise discrimination was quantified using average precision (AP), which summarises how well scores rank linked pairs above unlinked pairs by averaging precision over increases in recall. A pair was labelled positive if the two cases were connected by direct transmission or were sibling infectees sharing the same parent in the transmission tree. Positive pairs comprised approximately 0.14\% of all sampled pairs, making precision-recall summaries informative in this highly imbalanced setting \citep{Saito2015}. The pair labels are constructed separately from the BCubed reference memberships.

Temporal stability was evaluated by grouping cases by the week in which they became available and recomputing clustering after each cumulative weekly update. Each update included all cases available by that time. For every case present in both updates, we compared the earlier and later clusters containing that same case. Forward overlap is their shared membership divided by the later cluster's size; backward overlap uses the earlier cluster's size. Jaccard overlap divides shared membership by the number of distinct cases in the union of the two clusters. We averaged each measure equally across the shared cases, retaining newly added cases in the later-cluster denominators. This is a comparison anchored on each shared case, rather than a search for a best-matching cluster. The reported temporal summary then averages Jaccard overlap across weekly transitions. For the logistic comparators, training labels were available only for the first three weekly updates, comprising 86 initial cases. Each model's resolution was selected using BCubed F1 against the corrected reference memberships on these initial cases.

\subsection{Application to empirical dataset from Boston}\label{sec:ch3_epilink_boston}

The empirical evaluation used 772 high-quality SARS-CoV-2 whole-genome sequences ($>98\%$ coverage) collected by Massachusetts General Hospital (MGH)\nomenclature{MGH}{Massachusetts General Hospital} and the Massachusetts Department of Public Health (DPH)\nomenclature{DPH}{Massachusetts Department of Public Health} during the first wave of the COVID-19 pandemic in Boston, Massachusetts (March to May 2020) \citep{Lemieux2021}. Sequences were aligned using Nextclade against the SARS-CoV-2 reference genome, GenBank accession \path|MN908947.3| \citep{Aksamentov2021}.

Pairwise genetic distances were computed with \texttt{tn93}, a standalone software implementation of the Tamura-Nei 1993 (TN93) substitution model \citep{TamuraNei1993} incorporated in the HIV Transmission Cluster Engine (\texttt{HIV-TRACE})\nomenclature{HIV-TRACE}{HIV Transmission Cluster Engine} \citep{KosakovskyPond2018}. Ambiguous nucleotide codes represent more than one possible base at an alignment position. With \texttt{-a resolve}, compatible ambiguities were resolved to minimise the pairwise distance where possible, and the maximum tolerated fraction of sites containing resolvable ambiguities was set to 0.05 (\texttt{-g 0.05}). Only comparisons with at least 500 nucleotides of overlapping sequence were processed (\texttt{-l 500}), and distances were reported only below 0.0005 substitutions per site (\texttt{-t 0.0005}). TN93 distances were expressed as approximate numbers of nucleotide differences by multiplying by the alignment length of 29,903 nucleotides. These scaled distances are not directly counted single-nucleotide polymorphisms (SNPs)\nomenclature{SNP}{single-nucleotide polymorphism}. The temporal input for Boston was the absolute difference between collection dates. This orders a pair by sampling time for scoring and does not establish transmission direction.

Sequence metadata included collection date and exposure labels derived from epidemiological investigation, including Conference, skilled nursing facility (SNF)\nomenclature{SNF}{skilled nursing facility}, Boston Health Care for the Homeless Program (BHCHP)\nomenclature{BHCHP}{Boston Health Care for the Homeless Program}, City, and Unlabelled.

EpiLink compatibility scores were computed for the 276,056 unique pairs reported by \texttt{tn93} using the stochastic variant under baseline inference parameters. A weighted graph was constructed with minimum compatibility edge weight $10^{-4}$, and the Leiden algorithm was run at resolution $\gamma=0.2$. This resolution minimised the mean shortfall from each model's best BCubed F1 across the six models in the corrected synthetic evaluation, using the 86 cases available in the first three weekly updates. Candidate resolutions ranged from 0.1 to 1.0 in steps of 0.1, with ties resolved in favour of the smaller value (Figure~\ref{fig:resolution_regret}).

\begin{figure}[htbp]
    \centering
    \includegraphics[width=\textwidth]{epilink/resolution_regret}
    \caption[Difference from the best BCubed F1 across Leiden resolution values]{\textbf{Difference from the best BCubed F1 across Leiden resolution values.} For each of the six model variants in Table~\ref{tab:epilink_variants}, Leiden partitions were recomputed across candidate resolution values and compared with overlapping reference neighbourhoods derived from the transmission tree. The plotted difference is the model-specific maximum BCubed F1 minus the F1 achieved at each candidate resolution; it is not normalised. The solid line shows the mean difference across models, the shaded band the interquartile range, and the vertical dashed line the selected resolution of $\gamma=0.2$. The sweep used the 86 initial cases and corrected reference memberships. This value minimised the mean difference from the model-specific maxima and was used as the fixed resolution for the Boston analyses.}\label{fig:resolution_regret}
\end{figure}

To complement the EpiLink-Leiden analysis, we also evaluated TreeCluster on the Boston dataset. A maximum-likelihood tree was inferred from the aligned genomes using IQ-TREE 3 \citep{Wong2026}, then calibrated with TreeTime using collection dates. TreeCluster was applied to the dated tree using the \texttt{avg\_clade} and \texttt{max\_clade} methods, which respectively assign a cluster if the average or maximum pairwise distance within a clade is below a specified threshold.

Clusters of size $\ge 2$ were retained. For each cluster containing at least one Conference- or SNF-linked case, a $\chi^2$ goodness-of-fit test compared its exposure-category counts with the proportions among all 772 sequences, including the cluster itself. Genetic-distance and sampling-gap summaries were calculated from retained graph edges, either within a cluster or connecting it to sequences outside that cluster.

\subsection{Software and data availability statement}\label{sec:ch3_software_data}

The EpiLink software is available under an MIT\nomenclature{MIT}{Massachusetts Institute of Technology (licence name)} licence and archived on Zenodo with digital object identifier (DOI) \path|10.5281/zenodo.20402076|. Code, configuration files, and derived outputs required to reproduce the analyses are available in the evaluation branch and archived on Zenodo (DOI \path|10.5281/zenodo.20546213|). The Boston empirical sequences and metadata were derived from the previously published dataset \citep{Lemieux2021}. The synthetic benchmark was derived from SCoVMod outputs described by \citet{Banks2022}.

\subsection{Ethics statement}\label{sec:ch3_ethics}

The SARS-CoV-2 genomic sequences and epidemiological metadata analysed here were collected and anonymised by Massachusetts General Hospital and the Massachusetts Department of Public Health as part of the study reported by \citet{Lemieux2021}, which received approval from the MGH Institutional Review Board. The present work constitutes a secondary analysis of that published dataset and did not involve the collection of new human data or biological specimens; no additional ethics approval was required.

\section{Results}\label{sec:ch3_results}

\subsection{EpiLink produces interpretable compatibility surfaces}\label{sec:ch3_compatibility_surfaces}

Compatibility scores were highest among short sampling gaps and low-to-moderate genetic distances, with scores falling for pairs far from the time and genetic ranges implied by the selected scenarios (Figure~\ref{fig:epilink_surfaces}). The plots showed two broad regions with high scores: one at very short time gaps and modest genetic distance, consistent with co-primary infection, and one at slightly longer gaps but minimal genetic distance, consistent with direct transmission. The stochastic formulation (Figure~\ref{fig:epilink_surfaces}B) produced broader, less sharply bounded regions than the deterministic formulation (Figure~\ref{fig:epilink_surfaces}A). EpiLink therefore assigns the highest compatibility scores to pairs whose observed differences fall within the central ranges of either scenario, and much lower scores to pairs outside those ranges.

\begin{figure}[htbp]
    \centering
    \includegraphics[width=\textwidth]{surfaces}
    \caption[EpiLink compatibility surfaces]{\textbf{EpiLink compatibility surfaces.} Heatmaps of EpiLink compatibility under baseline inference parameters, evaluated at genetic distances and sampling-time differences. Each square cell is centred on a pair and spans one SNP by one day. Both panels use the same colour scale. (A) deterministic genetic component; (B) stochastic genetic component.}\label{fig:epilink_surfaces}
\end{figure}

\subsection{Baseline performance and computational cost}\label{sec:ch3_baseline_performance}

The synthetic evaluation used a 4,990-node transmission tree yielding 12,447,555 scored pairs, of which approximately 0.145\% were positive. At the selected sparsification threshold of $10^{-4}$, 87--97\% of edges were removed while more than 99.8\% of total graph weight was preserved. Recorded graph-construction and clustering time fell from 18.2--26.8 seconds to 0.68--2.60 seconds per model (10--39$\times$ speed-up). We retained this fixed threshold for the subsequent synthetic and Boston analyses. In the reproduced Boston graph it retained 99.99\% of the total score weight among the 276,056 reported pairs. This supports its limited numerical impact on total weight in these datasets, but does not establish an optimal threshold or demonstrate stability of every cluster.

In the baseline scenario, all six models substantially exceeded chance discrimination, with relative APs ranging from $63\times$ to $393\times$ (Table~\ref{tab:epilink_baseline}; Figure~\ref{fig:epilink_pr}). LD achieved the highest pairwise discrimination and the highest BCubed F1 (0.688). ESD was the strongest EpiLink variant on both measures, with F1 = 0.601, a difference of 0.087 from LD. Each deterministic-input variant outperformed its corresponding stochastic-input variant in both discrimination and clustering. The stochastic-input variants achieved BCubed F1 scores of 0.476--0.570.

As a supplementary comparison with clustering of a dated tree, TreeCluster achieved F1 = 0.665 on deterministic synthetic data (\texttt{max\_clade}, 35 days), below LD but above ESD and EDD. On stochastic synthetic data, TreeCluster achieved F1 = 0.615 (\texttt{avg\_clade}, 28 days), exceeding LS, ESS, and EDS. These comparisons use each method's best result over its candidate settings. Because TreeCluster assigns each sequence to one cluster without returning pairwise scores, AP was not computed. Additional TreeCluster comparisons are provided in Appendix Section~\ref{sec:app_ch3_supplementary_evaluation}.

\begin{table}[htbp]
    \centering
    \caption[Baseline performance metrics for all six models]{Baseline performance metrics for all six models. AP is average precision with 95\% confidence interval; F1 is BCubed F1; relative AP is AP divided by the prevalence of positive pairs (0.00145); resolution stability is mean F1 between cluster assignments at consecutive Leiden resolution values; temporal stability averages the weekly means of per-case Jaccard overlap. Models are ordered by AP.}\label{tab:epilink_baseline} % chktex 13
    \begin{thesistablebody}{@{}lccccc@{}}
    \toprule
    \textbf{Model} & \textbf{AP (95\% CI)} & \textbf{Relative AP} & \textbf{F1} & \textbf{Resolution stability} & \textbf{Temporal stability} \\
    \midrule
    LD & 0.568 (0.560--0.576) & 393 & 0.688 & 0.896 & 0.887 \\
    ESD & 0.469 (0.462--0.477) & 324 & 0.601 & 0.851 & 0.851 \\
    EDD & 0.300 (0.295--0.305) & 207 & 0.556 & 0.871 & 0.876 \\
    LS & 0.275 (0.271--0.281) & 190 & 0.570 & 0.927 & 0.926 \\
    ESS & 0.190 (0.187--0.194) & 132 & 0.534 & 0.830 & 0.801 \\
    EDS & 0.090 (0.089--0.092) & 63 & 0.476 & 0.834 & 0.755 \\
    \bottomrule
    \end{thesistablebody}
\end{table}

\begin{figure}[htbp]
    \centering
    \includegraphics[width=\textwidth]{baseline}
    \caption[Baseline synthetic performance across all six models]{\textbf{Baseline synthetic performance across all six models.} Precision-recall analysis assessing how well pairwise scores rank linked pairs above unlinked pairs, using EpiLink scores and logistic-regression probabilities. Performance is summarised with average precision and 95\% confidence interval.}\label{fig:epilink_pr}
\end{figure}

\subsection{Sensitivity to epidemiological and evolutionary assumptions}\label{sec:ch3_resolution_limits}

Averaged across the twelve perturbation scenarios, LD retained the highest AP and F1 under both inference conditions, and ESD retained the highest values among the EpiLink variants. Reduced substitution rate and shortened incubation mean reduced F1 in all six models under both conditions (Figure~\ref{fig:epilink_sensitivity}). For the EpiLink variants, losses under these two perturbations were larger with mismatched parameters. The largest F1 losses occurred when incubation CV was increased under mismatch.

\begin{figure}[htbp]
    \centering
    \includegraphics[width=\textwidth]{f1_loss}
    \caption[Sensitivity of F1 performance to perturbations]{\textbf{Sensitivity of F1 performance to perturbations of data-generating parameters.} Relative change in the best F1 score compared with baseline, calculated as (scenario F1 minus baseline F1) divided by baseline F1, for each model across twelve perturbation scenarios. In the matched condition, inference parameters were updated to reflect each scenario; in the mismatched condition, they were held fixed at baseline. Positive values indicate improvement; negative values indicate reduced performance.}\label{fig:epilink_sensitivity}
\end{figure}

Increased incubation CV produced sharply different responses under the two inference conditions. Under matched conditions, EDS gained 25.9\% in F1 and 253\% in AP, corresponding to absolute gains of 0.123 and 0.229, while LD lost 8.0\% in F1. Under mismatch, all six models suffered relative F1 losses of 55.7--65.6\%, with best F1 values of 0.202--0.237. At these selected settings, precision remained high (0.979--0.997), but recall was only 0.112--0.134. Thus the corrected reference revealed poor recovery of the reference groupings on both deterministic and stochastic inputs, even where pairwise AP improved.

Testing-delay mean and CV perturbations produced modest relative F1 changes, ranging from -3.7\% to +4.5\% across both conditions.

\subsection{Boston clusters align with known outbreak groups}\label{sec:ch3_boston_results}

Applying EpiLink to 276,056 unique pairs reported by \texttt{tn93} yielded 68 clusters of size $\ge 2$, containing 661 sequences in total. Most clusters were small, with a few much larger clusters (median = 3, mean = 9.7, max = 96) (Figure~\ref{fig:epilink_boston}A). The remaining 111 sequences were singletons. All four clusters containing Conference- or SNF-linked cases differed significantly from the exposure-category distribution among all 772 sequences ($\chi^2$, all $p<10^{-4}$) (Figure~\ref{fig:epilink_boston}B).

\begin{figure}[htbp]
    \centering
    \includegraphics[width=\textwidth]{boston}
    \caption[EpiLink clustering of the Boston MGH/DPH dataset]{\textbf{EpiLink clustering of the Boston MGH/DPH SARS-CoV-2 dataset.} (A) Size distribution of the 68 Leiden clusters containing at least two sequences (661 of the 772 sequences); four contained at least one Conference- or SNF-linked case. The remaining 111 sequences were singletons and are excluded from this panel. (B) Exposure-category counts and totals for the four clusters.}\label{fig:epilink_boston}
\end{figure}

The two largest clusters containing known outbreak-associated cases illustrate the signal EpiLink recovers. Cluster 7 ($n=96$) was dominated by SNF-linked cases (78.1\%). Retained edges within this cluster had a mean genetic distance of approximately 1.07 nucleotide differences (max 6) and a mean sampling gap of 2.81 days (max 14), consistent with an outbreak whose sampled cases were genetically similar and collected within a short period. Cluster 21 ($n=31$) contained 87.1\% Conference-linked cases and had similarly small genetic distances and short sampling gaps (mean genetic distance 1.26 nucleotide differences, max 5; mean gap 1.94 days, max 10). These genetic distances are scaled TN93 estimates. For both clusters, retained edges to sequences outside the cluster had a mean genetic distance greater than three nucleotide differences, although some outside-cluster matches had zero distance. Cluster membership was therefore associated with documented exposure groups; this association does not rule out an influence of sampling on the clusters.

The supplementary TreeCluster analysis of the Boston dataset showed a similar but more fragmented exposure-associated signal. TreeCluster with \texttt{avg\_clade} at 28 days yielded 111 clusters of size $\ge 2$, containing 542 sequences, with 230 singletons. Fifteen TreeCluster clusters contained at least one Conference- or SNF-linked case. The largest SNF-associated cluster contained 71 sequences, including 69 SNF-linked cases, while the largest Conference-associated cluster contained 12 sequences, including 7 Conference-linked cases. Compared with EpiLink, TreeCluster assigned fewer samples to non-singleton clusters and split the exposure-associated signal across more clusters. Detailed partition-agreement results are provided in Appendix Section~\ref{sec:app_ch3_supplementary_evaluation}. The largest SNF-associated groups shared 69 sequences: 69/96 (71.9\%) of EpiLink cluster 7 and 69/71 (97.2\%) of the TreeCluster group, with Jaccard overlap 69/98 = 0.704. The largest conference-associated groups shared seven sequences: 7/31 (22.6\%) of EpiLink cluster 21 and 7/12 (58.3\%) of the TreeCluster group, with Jaccard overlap 7/36 = 0.194. These percentages compare sample membership between methods; they differ from the exposure proportions within each cluster. Appendix Table~\ref{tab:app_ch3_boston_overlap} gives the greatest-overlap TreeCluster match for each of the four EpiLink focus clusters. EpiLink memberships were regenerated at the corrected selected resolution and checked against all regenerated cluster sizes and focus-cluster summaries before calculating overlap.

\section{Discussion}\label{sec:ch3_discussion}

EpiLink calculates graded recent-transmission compatibility scores without labelled training pairs, using explicit biological assumptions. Graph construction still requires an edge-weight threshold and a Leiden resolution, and the Boston resolution was selected using the labelled synthetic benchmark. The results show where this approach is useful and where its interpretation remains limited.

EpiLink was motivated partly by the challenge of superspreading events, where rapid transmission can produce many genetically similar cases within a short period. Sequence similarity alone may then be insufficient to distinguish recent transmission from coincidental similarity. In the Boston dataset, EpiLink grouped many cases with documented skilled nursing facility and conference exposures, and cluster membership was associated with those exposure groups \citep{Lemieux2021}. This illustrates its intended use when consensus genomes and case times are available but transmission labels are absent: identifying groups compatible with recent transmission, without reconstructing who infected whom.

The synthetic benchmark places that result in context. LD achieved the strongest pairwise discrimination and a BCubed F1 of 0.688, compared with 0.601 for ESD, the strongest EpiLink variant. The supervised comparator therefore achieved higher BCubed F1 in this benchmark. ESD uses stochastic-genetics inference on deterministic genetic inputs, so its performance does not imply that all data-generation and inference assumptions matched. The comparison also includes logistic training pairs in evaluation, and resolutions were selected using the reference memberships; it is not a held-out estimate of generalisation.

TreeCluster provided a comparison with clustering based on dated-tree distance thresholds. On deterministic inputs, its best F1 lay between LD and the EpiLink variants; on stochastic inputs, it exceeded both the logistic and EpiLink comparators. Thus the synthetic evaluation does not establish that compatibility-weighted graph clustering is more accurate than dated-tree clustering. In Boston, TreeCluster recovered exposure-associated groups but split them into more, smaller clusters than EpiLink. The large SNF groups overlapped substantially, whereas agreement between the largest conference-associated groups was lower. These results support shared recovery of some exposure-associated membership, without implying that both approaches identify equivalent groups. EpiLink additionally supplied a graded pairwise score describing compatibility with specified recent-transmission scenarios.% chktex 13

The corrected temporal evaluation gave mean Jaccard overlap of 0.755--0.926 across the six models. EDD had the highest mean among EpiLink variants (0.876), while ESD had the highest minimum (0.412). The largest changes generally occurred in the earliest updates, but EDS and ESS also showed appreciable turnover later. LS was the only model whose weekly mean Jaccard overlap remained above 0.7 across all 25 transitions. Differences between forward and backward overlap are consistent with new cases joining existing clusters as the dataset grows. This evaluation retained earlier cases at every update; it did not test a rolling window in which older cases leave the analysis.

The sensitivity analysis identified reduced substitution rate and shortened incubation mean as the changes that most consistently reduced performance. A lower substitution rate reduces the mutations expected over a given interval, while a shorter incubation period reduces the time available for mutation accumulation before sampling in these simulations. Both reduce the genetic information available to distinguish transmission relationships. More broadly, limited genetic variation constrains how confidently transmission relationships can be distinguished from sequences alone \citep{Louca2021,Farjo2023,Bendall2022}. Losses were generally larger under parameter mismatch. Performance also declined when inference used the data-generating parameter values, showing that mismatch was not the only source of reduced performance in these simulations.

Incubation-CV mismatch caused large clustering losses on both deterministic and stochastic inputs. The low recall despite high precision contrasts with the AP improvements retained by several models, showing that improved pairwise ranking did not ensure recovery of the reference neighbourhoods. This supports interpreting performance in light of the assumptions used to generate and analyse the data \citep{Campbell2018a,VillabonaArenas2020,Duchene2020,Suster2022}. The D/S distinction concerns how synthetic data were generated; it does not identify an inference strategy that an analyst can choose for observed genomes. Comparing how ED and ES handle uncertainty requires evaluating them on the same inputs.

Removing rate variation from data generation had different effects across inference conditions and variants. With inference matched to a strict clock, F1 increased by 2.8--5.8\% for EDD, ESD, and ESS, but decreased by 9.0\% for EDS. When inference retained the baseline relaxed clock, all four EpiLink variants gained F1, by 0.7--4.2\%. These results describe the tested simulations and do not establish that modelling rate variation is unnecessary for real data. Testing-delay perturbations had smaller effects on F1 over the range examined.

Several limitations remain. EpiLink operates on consensus genomes and therefore omits within-host diversity and minority variants that may carry additional transmission information. The user must specify the target scenario set $\mathcal{S}_\star$; the selection used here covers direct transmission and co-primary infection but does not explicitly score scenarios with unsampled intermediates. Compatibility scores are not calibrated posterior probabilities, and clustering depends on the Leiden resolution. Across the tested resolutions of 0.1--1.0, mean F1 agreement between consecutive cluster assignments ranged from 0.830 to 0.927 across models. This measures sensitivity to resolution choice, rather than establishing that any particular resolution is correct. Finally, the synthetic benchmark used a fixed transmission tree. Performance may differ for trees with different branching patterns or levels of variation in onward transmission.

EpiLink is useful when labelled training pairs are unavailable and the aim is to identify pairs compatible with recent transmission, including in outbreaks where some cases cause many onward infections. The score represents variation in infection timing, testing delays, and mutation accumulation under explicit parameter settings. Sensitivity analysis examines some consequences of changing those settings, without resolving all uncertainty about the model. This gives analysts a basis for explaining why a pair receives a high compatibility score and for assessing how that score changes under alternative assumptions.

\section{Conclusion}\label{sec:ch3_conclusions}

EpiLink combines genomic differences and sampling times into graded compatibility scores for defined recent linkage scenarios, without fitting to labelled transmission pairs. These scores remain conditional on the selected scenarios and model assumptions; they are not probabilities of transmission and do not identify who infected whom.

The corrected baseline evaluation places ESD highest among the EpiLink variants in BCubed F1, but below the supervised logistic comparator. The comparator also distinguished linked pairs more accurately in the separate pairwise analysis. Dated-tree TreeCluster achieved higher F1 than the EpiLink variants on both input types. In the Boston dataset, EpiLink recovered clusters enriched for documented conference and skilled nursing facility outbreaks. Together, these findings support using the scores to recover epidemiologically informative groups of infections, while limiting claims of comparative clustering accuracy.

Performance declined with reduced genetic resolution, and some parameter mismatches caused substantial losses in accuracy. The synthetic benchmark used a fixed transmission tree and a shared process-based model, so it does not establish performance across all outbreak settings.

Chapter~\ref{ch:genomic_networks} applies this compatibility measure to Scottish surveillance data and examines the resulting graphs under changing sampling conditions and viral lineages.
